\section{Conclusions and future work}\label{sec:conclu}

We presented a MBSE workflow that streamlines power system stabilizer design and deployment on low-cost microcontrollers leveraging the eFMI Standard. By introducing the \texttt{OpenIPSL\_CHIL} library, we established a ``single source of truth'' from physics-based, acausal Modelica models to production C code, clearly separating physics modeling from embedded development. All models, eFMUs, and C projects are provided open-source\footref{fnt:pss-repository} as an extensible workbench framework ultimately empowering researchers to rapidly prototype grid controllers with guaranteed traceability and significantly reduced capital expenditure.

Our systematic MiL, SiL, and CHiL design processes do not only guide development from off-line, continuous simulation models toward safety-critical, hard real-time suited, sampled embedded application, but also enable validation of intermediate designs of the various abstraction levels, significantly reducing the risk of late-stage issues with hard to trace sources of error. Verifying functional correctness via SiL testing provides strong confidence for deployment; once SiL tests pass, HiL system integration mainly involves checking memory fit and real-time execution on the target.

The mostly automatized MBSE process thereby supports iterative development cycles from Modelica models to deployable embedded code, enabling rapid prototyping for target switches and remodeling. The presented plant, for example, had to be redeployed on several targets due to real-time violations and IO pin restrictions before finally settling with the NUCLEO-H723ZG; and likewise the controller design underwent several iterations for signal-noise handling and real-time suited parameter-fitting before CHiL results have been satisfactory. And instead of using the built-in ADCs and DACs of the selected microcontrollers as well as the power supplies provided onboard, we could improve signal-to-noise ratios by using dedicated external ADC and DAC components with careful design without further microcontroller upgrades; such hardware setup redesign would have to be reflected in the model (e.g., filter configurations, parameters, etc), confirming the need of an MBSE-supported iterative development process along the lines of eFMI.

Besides systematic functional testing, we also validated the non-functional quality of the production code generated by the used eFMI tooling. Production code analyses confirm MISRA C:2023 and SEI CERT C Coding Standard\footref{fnt:sei-cert-c} compliance; execution time measurements confirm real-time capability; and the rigorous error signaling of eFMI GALEC confirms the absence of numerical issues.

Last but not least, our approach dramatically reduces capital expenditure. It reduces software licenses to only Dymola (with \textit{Dymola Source Code Generation} option) while using free development tools for the STM32 microcontrollers. And although eFMI was not designed for real-time plant simulation, we show that it has strong potential for that use, eventually replacing proprietary real-time platforms that typically cost $\gg$\$50,000 with $<$\$50 microcontrollers. This came as a surprise to us. Our original plan was to use one of the established real-time platforms, like from dSPACE, for the power generator plant-model (and only deploy the PSS controller-model on a microcontroller); we only tried eFMI-based deployment of the plant-model out of curiosity how far we can get with Dymola's eFMI tooling, but soon figured this may indeed work out!

Future research will concentrate on both the modeling and eFMU integration sides. We intend to work on scaling our current proof-of-concept to cover a broader spectrum of PSS models, and on creating templates that support a variety of embedded hardware boards from the STM32 family. An important objective is to exploit the eFMI recalibration to enable online parameter tuning, thereby allowing the PSS to adapt its behavior to real-time grid conditions. Finally, we plan to showcase the integration of eFMUs (exported as source code FMUs or the eFMI production code directly system integrated) into industrial real-time platforms from dSPACE, Typhoon HIL, and OPAL-RT Technologies.
